%% =================================================================
%% Beber et al. IEEE Robotics and Automation Letters, Feb 2026.
%% Reconstruction of page 2 (printed p. 2).
%% Contents: continuation of Section I (four-bullet contributions
%% list), Section II Related Work (II-A Bayesian Optimisation for
%% Palpation, II-B Ergodic Exploration), start of Section III
%% Methodology.
%% =================================================================

\documentclass[11pt]{article}

\usepackage[margin=0.85in]{geometry}
\usepackage{amsmath, amssymb}
\usepackage{mathrsfs}
\usepackage{xcolor}
\usepackage{fancyhdr}

\pagestyle{fancy}
\fancyhf{}
\lhead{\small Beber et al.}
\rhead{\small IEEE RAL, Accepted Feb 2026}
\cfoot{\small 2 $\to$ \arabic{page}}
\renewcommand{\headrulewidth}{0.4pt}

\setcounter{page}{1}
\setcounter{equation}{0}
\setlength{\parskip}{6pt}
\setlength{\parindent}{0pt}

\begin{document}

% ---------- Continuation of Section I Introduction ----------
shelf F/T sensor ensures portability, cost-effectiveness, and broader
applicability of the proposed framework. We couple a force-based
contact model implemented via the Dimensionality Reduction Method
(DRM) [18] with an ergodic exploration strategy to efficiently
construct stiffness maps, where the robot explores the environment
in proportion to the Expected Information Density (EID). To plan
trajectories, we employ the Heat Equation Driven Area Coverage
(HEDAC) algorithm [19], a real-time ergodic controller that balances
exploration and exploitation at 100 Hz. This replaces earlier ergodic
control methods, like Spectral Multiscale Coverage (SMC), which was
limited to 4 Hz as reported in [16]. The proposed framework is
evaluated in simulation and real-world experiments, including
validation on a silicone sample with a stiff inclusion (Fig.~1),
demonstrating robust performance and, in some cases, superior
detection and boundary delineation compared to BO.

To summarise, the main contributions of this paper are:
\begin{itemize}
\item The design of an EID tailored to autonomous stiffness mapping,
  which explicitly combines uncertainty, stiffness magnitude, and
  spatial gradients to guide ergodic exploration toward diagnostically
  relevant regions;
\item A closed-loop combination between online force-based
  viscoelastic parameter estimation and ergodic trajectory planning,
  where stiffness estimates continuously reshape the information
  objective driving exploration;
\item A real-time implementation of ergodic palpation using HEDAC,
  enabling continuous exploration--exploitation trade-offs at 100 Hz
  with continuous motion;
\item A convergence-based stopping criterion derived from the
  ergodic metric.
\end{itemize}

% ---------- Section II ----------
\section*{II.\ Related Work}

Typically, search algorithms rely on three key components: (i) a
model of the underlying distribution, (ii) a strategy to select the
next sampling location, and (iii) a mechanism to store and update the
collected data. Although Gaussian Process Regression (GPR) [20] is
widely adopted to model spatially varying physical properties (ii),
such as tissue stiffness, the choice of sampling strategy (i) and
planning (iii) remains an active area of research.

\subsection*{A.\ Bayesian Optimisation for Palpation}

Several studies have employed BO to guide robotic palpation to reduce
the number of probing actions while maximising information gain. For
example, Garg et al.\ [12] investigated how different acquisition
functions affect the estimated stiffness map, showing how their
choice influences both peak detection and regional estimation. Yan
et al.\ [15] highlighted a key limitation of BO: while it excels at
locating stiffness peaks, it struggles to delineate their boundaries
accurately. To address this, they proposed a two-step approach
combining BO-based search and boundary refinement using a radial
strategy originating from the centroid of the estimated region.
However, this strategy requires a reliable estimation of the
centroid, which may not be guaranteed after exploration, and may be
less adaptive to irregular or non-convex boundaries. Expected
Improvement has been used as an acquisition function to improve
stiffness estimation [13], and later extended with a utility-guided
approach [16] that incorporates prior beliefs and penalises excessive
motion. Chalasani et al.\ proposed a method that employs GPR to
estimate tissue stiffness, which is then used as a prior to fit a
second GPR model for reconstructing the tissue surface geometry [21].
In a follow-up study, they extended this approach to support online
stiffness mapping during surgery [22]. A critical limitation of BO
is that exploration and exploitation are handled sequentially through
the choice of the acquisition function. As a result, BO may focus
excessively on promising regions, thereby overlooking smaller or
subtler areas that could yield informative measurements.

\subsection*{B.\ Ergodic Exploration}

Ergodic control has emerged as a compelling alternative to BO for
tasks that involve spatially distributed information, where
exploration must account for system dynamics. Unlike point-wise
optimisation strategies, which select discrete sampling locations,
ergodic control generates trajectories that allocate time across the
environment proportionally to the EID, making it well suited for
continuous spatial estimation tasks. Mathew and Mezi\'c [23]
introduced the Spectral Multiscale Coverage (SMC) framework, using
Fourier-based metrics to quantify how well a trajectory covers a
given spatial distribution. Miller et al.\ [24] demonstrated the use
of ergodic control for information-guided exploration, integrating
the Linear Quadratic Regulator (LQR) formulation of SMC [25].
However, the high computational cost of optimisation limited
real-time adaptation. Ayvali et al.\ [16] presented the first
application of ergodic control to palpation tasks using SMC,
concluding that BO achieved superior performance, particularly when
prior information was available. Conversely, our results indicate
the opposite trend, highlighting the potential of ergodic control
for robotic palpation. The main reason is that SMC often prioritises
global exploration and can suffer from over-sampling already explored
regions, while its reliance on spectral transforms poses challenges
for real-time deployment. To address these limitations, Ivic et al.\
[19] proposed HEDAC, which leverages radial basis functions and a
potential field derived from a stationary heat equation to compute
smooth local trajectories. This method improves scalability and
enables real-time implementation, avoiding the computational burden
of Fourier analysis.

% ---------- Section III opening ----------
\section*{III.\ Methodology}

In this work, we combine online force-based parameter estimation
using a viscoelastic tissue model with real-time ergodic control
implemented via HEDAC, enabling continuous robotic palpation without
tactile sensors and without sacrificing spatial coverage or
responsiveness. While ergodic control and online viscoelastic
estimation have been previously investigated, we show that effective
ergodic palpation critically depends on the definition of the target
distribution guiding exploration. We therefore propose an EID
tailored to stiffness mapping and lesion detection, which balances
exploration and exploitation by jointly considering uncertainty,
stiffness magnitude, and

\end{document}
